\documentclass[11pt]{article}
\usepackage[T1]{fontenc}
\usepackage{textcomp}
\usepackage[margin=0.75in]{geometry}
\usepackage{amsmath,amssymb,amsthm}
\usepackage{array,booktabs}
\usepackage{hyperref}
\setlength{\parindent}{0pt}
\newtheorem{theorem}{Theorem}
\theoremstyle{definition}
\newtheorem{definition}{Definition}
\title{Flow Proof Artifact: prop-29-alternate-angles-equal}
\author{Generated by \texttt{flow doc proof}}
\date{Flow Proof Book}
\begin{document}
\maketitle
A straight line falling on parallel straight lines makes alternate angles equal.\\[0.5em]
\textbf{Source.} euclid — Elements, Book I, Proposition 29\\[0.5em]
\bigskip
\noindent{\large \textbf{Derived fact 1.} \textit{A straight line falling on parallel straight lines makes alternate angles equal} \hfill the Euclidean plane \cdot Euclid Book I \cdot Proposition 29: alternate angles are equal when a transversal crosses parallels}\par
\smallskip\noindent\textit{Coordinate: the Euclidean plane \cdot Euclid Book I \cdot Proposition 29: alternate angles are equal when a transversal crosses parallels}\par
\smallskip\noindent\textit{Source: euclid — Elements, Book I, Proposition 29}\par
\smallskip\noindent\textit{Built on: proposition 14: lines with interior angles less than two right angles meet, for Book I of the Elements in on the Euclidean plane, proposition 15: vertical angles are equal, for Book I of the Elements in on the Euclidean plane}\par
\medskip
\noindent\textbf{Goal.} A straight line falling on parallel straight lines makes alternate angles equal\par
\noindent$\displaystyle \alpha = \beta$\par
\medskip
\noindent
\begin{tabular}{@{} >{\raggedright\arraybackslash}p{0.06\textwidth} >{\raggedright\arraybackslash}p{0.47\textwidth} >{\raggedleft\arraybackslash}p{0.41\textwidth} @{}}
\textbf{\#} & \textbf{Proof} & \textbf{Mathematics} \\
\midrule
\textbf{1.} & We prove that proposition 29: alternate angles are equal when a transversal crosses parallels for Book I of the Elements in the Euclidean plane. & \\[0.45em]
\textbf{2.} & Let alpha = alternate\_angle\_AGH. & \\[0.45em]
\textbf{3.} & Let beta = alternate\_angle\_GHD. & \\[0.45em]
\textbf{4.} & We invoke the derived fact: lines\_AB\_and\_CD\_are\_parallel. & \\[0.45em]
\textbf{5.} & We invoke the derived fact: transversal\_EF\_crosses\_parallels\_at\_G\_and\_H. & \\[0.45em]
\textbf{6.} & We invoke the axiom governing Book I of the Elements in the Euclidean plane: proposition 14: lines with interior angles less than two right angles meet, for Book I of the Elements in on the Euclidean plane. & \\[0.45em]
\textbf{7.} & We invoke the derived fact governing Book I of the Elements in the Euclidean plane: proposition 15: vertical angles are equal, for Book I of the Elements in on the Euclidean plane. & \\[0.45em]
\textbf{8.} & From step 2, step 3, step 4, step 5, step 6, and step 7, by the chain of results established in Book I (step 2, step 3, step 4, step 5, step 6, and step 7), we obtain angle $\alpha$ equals angle $\beta$. Hence proven. & $\displaystyle \alpha = \beta$ \\[0.45em]
\bottomrule
\end{tabular}
\label{thm:Geometry-Euclid Book I-proposition_29_alternate_angles_are_equal_when_a_transversal_crosses_parallels}
\par\medskip\hrule
\end{document}
